\Needspace{14\baselineskip}
\section{Trabalhos Relacionados}
\label{sec:trabalhos-relacionados}

% PRISMA Item 3: Rationale - posicionamento contra revisões anteriores

Esta seção posiciona a presente revisão contra duas classes de
trabalhos secundários: surveys sobre memória para agentes
baseados em LLMs em geral e revisões sistemáticas sobre LLMs
aplicados a Engenharia de Software.
A primeira classe é representada principalmente pelo levantamento
recente de Du et al.; a segunda agrupa cinco SLRs publicadas
entre 2024 e 2025.
Nenhuma das duas classes adota a persistência inter-sessão como
objeto primário, nem reporta análise de custo operacional por
tarefa, que são as dimensões críticas avaliadas nesta investigação.
A \autoref{tab:slrs-previas} sintetiza a comparação por foco
temático e cobertura dessas duas dimensões.

\begin{table}[htbp]
\centering
\caption{Comparação entre trabalhos secundários relacionados e
esta revisão. As colunas \emph{Persistência?} e \emph{Custo?}
indicam, respectivamente, se o trabalho aborda explicitamente a
persistência de conhecimento entre sessões e se reporta análise
de custo operacional por tarefa.}
\label{tab:slrs-previas}
\small
\begin{tabularx}{\textwidth}{l c X c c}
\toprule
\textbf{Trabalho} & \textbf{Ano} & \textbf{Foco principal} &
  \textbf{Persistência?} & \textbf{Custo?} \\
\midrule
Du et al.~\cite{du2026memory} & 2026 &
  Memória para agentes LLM em geral &
  Sim & Não \\
Hou et al.~\cite{hou2024llmse} & 2024 &
  LLMs para Engenharia de Software (395 estudos) &
  Não & Não \\
Zhang et al.~\cite{zhang2024unifying} & 2024 &
  Modelos de linguagem para código (NLP $\cap$ ES) &
  Não & Não \\
Jin et al.~\cite{jin2024llmagents} & 2024 &
  Evolução de LLMs para agentes em ES &
  Não & Não \\
Wang et al.~\cite{wang2024llmagents} & 2024 &
% audit:ignore-next-line — describes cited work's own framing
  Landscape de agentes autônomos para ES &
  Parcial & Não \\
Liu et al.~\cite{liu2024llmagents} & 2024 &
  Agentes LLM para Engenharia de Software (124 estudos) &
  Parcial & Não \\
\midrule
\textbf{Esta revisão} & 2026 &
  \textbf{Persistência de conhecimento em agentes de codificação
  baseados em IA (28 estudos)} &
  \textbf{Sim} & \textbf{Sim} \\
\bottomrule
\end{tabularx}
\end{table}

\subsection{Surveys sobre memória para agentes LLM}
\label{sec:rel-memoria}

\citeonline{du2026memory} apresentam o levantamento mais recente
sobre mecanismos de memória para agentes baseados em LLMs,
propondo uma taxonomia tridimensional (escopo temporal, substrato
representacional e política de controle) que esta revisão adota e
refina como ponto de partida da QS1.
O survey cataloga cinco famílias de mecanismos de memória:
compressão de contexto residente, stores aumentados por
recuperação, memória reflexiva auto-corretiva, contextos
hierárquicos virtuais e gestão de memória aprendida por política.
O escopo dos autores é explicitamente amplo, incluindo agentes
para tarefas gerais como navegação web, raciocínio matemático e
diálogo, sem restrição ao domínio de codificação.
Como consequência, a cobertura específica para Engenharia de
Software no survey de Du et al.\ limita-se a ChatDev e MetaGPT,
enquanto sistemas centrais para a presente investigação, como
MemCoder~\cite{deng2026memcoder},
CodeMEM~\cite{wang2026codemem} e a abordagem de Repository
Memory de~\citeonline{wang2026improving}, não são mencionados.
A presente revisão complementa esse trabalho restringindo o
escopo a agentes de codificação e incorporando análise de custo
operacional por tarefa, dimensão ausente no survey de Du et al.

\subsection{Revisões sistemáticas sobre LLMs em Engenharia de
Software}
\label{sec:rel-llmse}

Cinco revisões sistemáticas recentes mapeiam aplicações de LLMs em
Engenharia de Software com escopo amplo.
\citeonline{hou2024llmse} cobrem 395 estudos publicados até 2023,
organizados por tarefa de engenharia (geração de código, reparo,
teste, manutenção) e por tipo de modelo, sem subseção dedicada a
memória ou persistência cross-session.
\citeonline{zhang2024unifying} unificam perspectivas de
Processamento de Linguagem Natural e Engenharia de Software sobre
modelos para código, com ênfase em pré-treinamento e benchmarks de
avaliação.
\citeonline{jin2024llmagents} traçam a transição de LLMs para
agentes baseados em LLMs em Engenharia de Software, mas tratam
memória como componente arquitetural genérico, sem caracterizar
arquiteturas de persistência inter-sessão.
\citeonline{wang2024llmagents} mapeiam o cenário dos agentes
autônomos para Engenharia de Software com foco em pipelines de
execução, e \citeonline{liu2024llmagents} cobrem 124 estudos com
discussão parcial sobre módulos de memória, porém sem síntese
cross-system de eficácia e custo da persistência.
Nenhuma dessas revisões responde diretamente às questões
secundárias QS1--QS5 desta investigação: a taxonomia de
arquiteturas de persistência, a comparabilidade de benchmarks
específicos para memória, a quantificação de ganhos sobre
baselines sem memória, a relação custo-eficácia ou os mecanismos
de degradação por memória.

\subsection{Lacuna posicional desta revisão}
\label{sec:rel-lacuna}

Esta revisão restringe o escopo aos agentes de codificação
baseados em IA e adota a persistência de conhecimento entre
sessões como objeto primário de investigação, não como componente
arquitetural acessório de uma síntese mais ampla.
Sobre essa base, incorpora a análise de custo operacional por
tarefa como dimensão dedicada (QS4), o que constitui o
diferencial central frente aos trabalhos anteriores: nenhuma das
SLRs prévias identificadas reúne simultaneamente foco em
persistência inter-sessão, escopo de agentes de codificação e
quantificação cross-study de custos.
Essa configuração permite responder questões que os trabalhos
secundários anteriores não abordam, sobretudo a ausência de
comparações head-to-head entre arquiteturas, a subnotificação
sistemática de custos e os mecanismos pelos quais a memória pode
prejudicar o desempenho do agente.
